\section{\texorpdfstring{$\mathcal{I}$}{I} 算子}\label{sec:vanishing-ideal-operator}

在第~\ref{sec:vanishing-set-operator}~节, 我们学习了怎样借助 $\mathcal{V}$ 算子, 给 $K[x_{1},\ldots,x_{n}]$ 的每个子集对应 $\mathbb{A}^{n}$ 的一个子集. 在本节中, 我们反过来描述一个算子, 给 $\mathbb{A}^{n}$ 的每个子集对应 $K[x_{1},\ldots,x_{n}]$ 的一个子集.

\begin{Definition}{消没理想}\leavevmode\label{def:vanishing-ideal}
  设 $X\subseteq\mathbb{A}^{n}$ 是一个子集. $X$ 的\textbf{消没理想}定义为
  \[
    \mathcal{I}(X)=\{f\in K[x_{1},\ldots,x_{n}]
      \mid f(a_{1},\ldots,a_{n})=0\text{ 对所有 }(a_{1},\ldots,a_{n})\in X\}.
  \]
\end{Definition}

换句话说, $X$ 的消没理想是所有在 $X$ 上处处为零的多项式所成的集合. 正如名称所暗示的那样, 集合 $\mathcal{I}(X)$ 不只是多项式环 $K[x_{1},\ldots,x_{n}]$ 的一个子集, 还是一个理想 (Exercise~\ref{ex:vanishing-set-of-polynomials-is-ideal}).

\begin{Example}{在 $(0,0)$ 处为零的多项式}\leavevmode\label{ex:vanishing-ideal-of-origin}
  设 $X=\{(0,0)\}\subseteq\mathbb{A}^{2}$. 多项式 $f\in K[x,y]$ 在 $(0,0)$ 处为零, 当且仅当 $f$ 的常数项为零. 如 Example~\ref{ex:vanishing-constant-ideal} 所述, 所有这样的多项式构成理想 $\langle x,y\rangle$. 因此, $\mathcal{I}(X)=\langle x,y\rangle$.
\end{Example}

\begin{Example}{一元情形中的消没理想}\leavevmode\label{ex:one-variable-vanishing-ideals}
  设 $X=\{1,3\}\subseteq\mathbb{A}^{1}_{\mathbb{R}}$. 根据 Corollary~\ref{cor:polynomial-factor-theorem}, 多项式 $f\in\mathbb{R}[x]$ 在 $1$ 处为零, 当且仅当 $x-1$ 整除 $f$; $f$ 在 $3$ 处为零, 当且仅当 $x-3$ 整除 $f$. 由于 $x-1$ 和 $x-3$ 都不可约, 由 $\mathbb{R}[x]$ 中分解的唯一性可得
  \[
    \mathcal{I}(X)=\{f\in\mathbb{R}[x]\mid(x-1)(x-3)\text{ 整除 }f\}
    =\langle x^{2}-4x+3\rangle.
  \]
  采用类似方法可求任意 $X\subseteq\mathbb{A}^{1}$ 的消没理想 (Exercise~\ref{ex:vanishing-ideals-of-subsets-of-line}).
\end{Example}

\begin{Example}{抛物线的消没理想}\leavevmode\label{ex:vanishing-ideal-of-parabola}
  设 $K$ 是无限域, 并设 $X$ 为下式定义的仿射簇:
  \[
    X=\mathcal{V}(y-x^{2})=\{(a,a^{2})\mid a\in K\}\subseteq\mathbb{A}^{2}.
  \]
  哪些多项式在 $X$ 的每个点都为零? 任何形如 $(y-x^{2})f(x,y)$ 的多项式当然在 $X$ 的每个点都为零, 所以 $\langle y-x^{2}\rangle\subseteq\mathcal{I}(X)$. 下面证明反向的包含关系也成立.

  设 $f\in\mathcal{I}(X)$. Example~\ref{ex:quotient-parabola-computations} 中的论证表明, 在商环 $K[x,y]/\langle y-x^{2}\rangle$ 中, 有 $[f(x,y)]=[f(x,x^{2})]$. 特别地, 这意味着
  \begin{equation}\label{eq:parabola-vanishing-ideal-difference}
    f(x,y)-f(x,x^{2})\in\langle y-x^{2}\rangle.
  \end{equation}
  由 $f$ 在 $X=\{(a,a^{2})\mid a\in K\}$ 上处处为零的假设可知, 对每个 $a\in K$ 都有 $f(a,a^{2})=0$. 这意味着 $f(x,x^{2})\in K[x]$ 是一个具有无穷多个零点的一元多项式. 只有在 $f(x,x^{2})$ 为零多项式时这才可能. 将 $f(x,x^{2})=0$ 代入 \eqref{eq:parabola-vanishing-ideal-difference}, 可知 $f\in\langle y-x^{2}\rangle$, 从而有 $\mathcal{I}(X)\subseteq\langle y-x^{2}\rangle$.

  两个方向的包含关系都已证明, 因此 $\mathcal{I}(X)=\langle y-x^{2}\rangle$.
\end{Example}

即使在 Example~\ref{ex:vanishing-ideal-of-parabola} 的抛物线这样简单的情形中, 证明消没理想等于 $\langle y-x^{2}\rangle$ 也已经颇费周折. 的确, 一般来说, 计算仿射簇的消没理想并不容易, 通常必须针对每个具体情形采用专门的方法 (更多此类计算的例子见 Exercises~\ref{ex:vanishing-ideals-of-subsets-of-line}--\ref{ex:unit-circle-infinite-field}). 如第~\ref{sec:nullstellensatz}~节所示, 零点定理会大大简化计算消没理想的工作.

现在我们已经介绍了 $\mathcal{V}$ 和 $\mathcal{I}$ 算子, 它们把 $K[x_{1},\ldots,x_{n}]$ 的子集与 $\mathbb{A}^{n}$ 的子集相互联系起来. 自然会问, 这两个算子在多大程度上互为逆运算? 下面的结果给出第一个答案: 在一般情形下, 它们不互为逆运算; 但把它们限制在消没理想和仿射簇上以后, 它们就互为逆运算了.

\begin{Proposition}{复合 $\mathcal{V}$ 和 $\mathcal{I}$ 算子}\leavevmode\label{prop:compose-vanishing-operators}
  \strut\par\nopagebreak[4]
  \begin{enumerate}[label=\arabic*.,ref=\arabic*]
    \item\label{item:compose-vanishing-operators-sets} 设 $X\subseteq\mathbb{A}^{n}$. 则
    \[
      \mathcal{V}(\mathcal{I}(X))\supseteq X,
    \]
    等号成立当且仅当对某个 $S\subseteq K[x_{1},\ldots,x_{n}]$, 有 $X=\mathcal{V}(S)$.
    \item\label{item:compose-vanishing-operators-polynomials} 设 $S\subseteq K[x_{1},\ldots,x_{n}]$. 则
    \[
      \mathcal{I}(\mathcal{V}(S))\supseteq S,
    \]
    等号成立当且仅当对某个 $X\subseteq\mathbb{A}^{n}$, 有 $S=\mathcal{I}(X)$.
  \end{enumerate}
\end{Proposition}

\begin{Proof}
  我们证明第~\ref{item:compose-vanishing-operators-sets}~部分, 第~\ref{item:compose-vanishing-operators-polynomials}~部分留给 Exercise~\ref{ex:prove-second-composition-inclusion}.

  取 $a=(a_{1},\ldots,a_{n})\in X$. 为证明 $a\in\mathcal{V}(\mathcal{I}(X))$, 任取多项式 $f\in\mathcal{I}(X)$. 根据消没理想的定义, 对所有 $b\in X$ 都有 $f(b)=0$. 特别地, $f(a)=0$. 这样便证明了对每个 $f\in\mathcal{I}(X)$ 都有 $f(a)=0$. 这意味着 $a\in\mathcal{V}(\mathcal{I}(X))$.

  若等号成立, 则 $X$ 是 $S=\mathcal{I}(X)$ 的零点集, 这就证明了充要条件的一个方向. 为了证明反向, 设对某个 $S\subseteq K[x_{1},\ldots,x_{n}]$, 有 $X=\mathcal{V}(S)$. 必须证明 $\mathcal{V}(\mathcal{I}(X))\subseteq X$. 这等价于证明 $a\notin X$ 蕴含 $a\notin\mathcal{V}(\mathcal{I}(X))$, 所以假定 $a\notin X$. 由于 $X=\mathcal{V}(S)$ 且 $a\notin X$, 存在某个 $f\in S$ 使得 $f(a)\ne0$. 又因为 $f\in S$ 且 $X=\mathcal{V}(S)$, $f$ 在 $X$ 上处处为零, 所以 $f\in\mathcal{I}(X)$. 由 $f\in\mathcal{I}(X)$ 且 $f(a)\ne0$, 可知 $a\notin\mathcal{V}(\mathcal{I}(X))$.
\end{Proof}

\BookPairedNoteLeft
  {若 $X\subseteq\mathbb{A}^{n}$ 是仿射簇, 该命题的第~\ref{item:compose-vanishing-operators-sets}~部分说明有一个以 $X$ 为零点集的特别的理想, 即 $\mathcal{I}(X)$.}
  {若 $X\subseteq\mathbb{A}^{n}$ 是一个子集, Proposition~\ref{prop:compose-vanishing-operators} 第一部分中的集合 $\mathcal{V}(\mathcal{I}(X))$ 还有另一种解释: 它是包含 $X$ 的最小仿射簇 (Exercise~\ref{ex:smallest-affine-variety-containing-set}). 从这个意义上说, 它类似于由集合 $S$ 生成的理想, 即包含 $S$ 的最小理想.}
\noindent 我们甚至可以称 $\mathcal{V}(\mathcal{I}(X))$ 为\textquotedbl{}由 $X$ 生成的仿射簇\textquotedbl{}, 并记作 $\langle X\rangle$. 用这个记号, Proposition~\ref{prop:compose-vanishing-operators} 的最后一部分就与 Proposition~\ref{prop:affine-varieties-defined-by-ideals} 相类似:
\[
  \mathcal{I}(X)=\mathcal{I}(\langle X\rangle).
\]

从 Proposition~\ref{prop:compose-vanishing-operators} 的两个充要条件可以看出, 这些包含关系当然可能是严格的. 例如, 若 $X$ 不是仿射簇, 则对任意 $S$ 都有 $X\ne\mathcal{V}(S)$, 因而 $X\ne\mathcal{V}(\mathcal{I}(X))$. 同样, 若 $S$ 不是理想, 就不可能有 $S=\mathcal{I}(\mathcal{V}(S))$, 因为后者是理想. 事实上, 即使 $S$ 是理想, 等号也不必成立. 下面的例子说明这一点.

\begin{Example}{$\mathcal{I}(\mathcal{V}(I))\ne I$}\leavevmode\label{ex:square-ideal-not-vanishing-ideal}
  设 $I=\langle x^{2}\rangle\subseteq K[x]$. 则
  \[
    \mathcal{V}(I)=\{a\in K\mid a^{2}=0\}.
  \]
  因为 $K$ 是域, 所以 $a^{2}=0$ 当且仅当 $a=0$. 因此, $\mathcal{V}(I)=\{0\}$. 与 Example~\ref{ex:one-variable-vanishing-ideals} 相同的推理表明 $\mathcal{I}(\{0\})=\langle x\rangle$, 所以
  \[
    \mathcal{I}(\mathcal{V}(I))=\mathcal{I}(\{0\})=\langle x\rangle.
  \]
  消没理想 $\mathcal{I}(\mathcal{V}(I))=\langle x\rangle$ 包含 $I=\langle x^{2}\rangle$, 但不等于它. 因此, 根据 Proposition~\ref{prop:compose-vanishing-operators} 第二部分中的充要条件, $I$ 不是任何 $X\subseteq\mathbb{A}^{1}$ 的消没理想.
\end{Example}

在第~\ref{sec:affine-varieties}~节末尾, 我们看到 $\mathcal{V}$ 算子给出满射
\[
  \mathcal{V}:\{K[x_{1},\ldots,x_{n}]\text{ 中的理想}\}
    \longrightarrow\{\mathbb{A}^{n}\text{ 中的仿射簇}\}.
\]
由 Proposition~\ref{prop:compose-vanishing-operators} 可知, 若把定义域限制为消没理想的集合, 即那些能作为某个集合的消没理想的理想, 那么 $\mathcal{V}$ 算子就成为双射, 其逆由 $\mathcal{I}$ 算子给出. 因此, 要理解代数与几何之间的这本词典, 我们的目标就是更好地理解 $K[x_{1},\ldots,x_{n}]$ 中的消没理想. 受 Example~\ref{ex:square-ideal-not-vanishing-ideal} 中观察的启发, 下一节将介绍根理想这一代数概念, 沿这个方向迈出第一步.

\BookExerciseHeading
\begin{BookExercises}
  \item\label{ex:vanishing-set-of-polynomials-is-ideal} 对任意子集 $X\subseteq\mathbb{A}^{n}$, 证明 $\mathcal{I}(X)\subseteq K[x_{1},\ldots,x_{n}]$ 是一个理想.
  \item\label{ex:prove-second-composition-inclusion} 证明 Proposition~\ref{prop:compose-vanishing-operators} 第~\ref{item:compose-vanishing-operators-polynomials}~部分.
  \item\label{ex:vanishing-ideals-of-subsets-of-line} 设 $X\subseteq\mathbb{A}^{1}$.
    \begin{enumerate}[label=(\alph*)]
      \item 假设 $X=\{a_{1},\ldots,a_{r}\}$ 是有限集. 利用 $K[x]$ 中分解的唯一性证明
        \[
          \mathcal{I}(X)=\langle(x-a_{1})\cdots(x-a_{r})\rangle\subseteq K[x].
        \]
      \item 假设 $X$ 是无限集. 证明 $\mathcal{I}(X)=\langle0\rangle$.
    \end{enumerate}
  \item\label{ex:vanishing-ideal-quadratic-real-complex} 分别求 $\mathcal{V}(x^{2}+1)\subseteq\mathbb{A}^{1}$ 的消没理想:
    \begin{enumerate}[label=(\alph*)]
      \item 在 $\mathbb{R}$ 上;
      \item 在 $\mathbb{C}$ 上.
    \end{enumerate}
  \item\label{ex:vanishing-ideal-of-a-single-point} 设 $X=\{(a_{1},\ldots,a_{n})\}\in\mathbb{A}^{n}$ 是一个单点. 证明
    \[
      \mathcal{I}(X)=\langle x_{1}-a_{1},\ldots,x_{n}-a_{n}\rangle.
    \]
  \item\label{ex:vanishing-ideal-of-real-unit-circle} 设 $X=\mathcal{V}(x^{2}+y^{2}-1)\subseteq\mathbb{A}^{2}_{\mathbb{R}}$ 是单位圆. 这道习题证明 $\mathcal{I}(X)=\langle x^{2}+y^{2}-1\rangle$.
    \begin{enumerate}[label=(\alph*)]
      \item 证明 $\mathcal{I}(X)\supseteq\langle x^{2}+y^{2}-1\rangle$.
      \item 证明 $\mathcal{I}(X)\subseteq\langle x^{2}+y^{2}-1\rangle$, 可以采用以下证明提纲. 假设 $f\in\mathcal{I}(X)$.
        \begin{enumerate}[label=\roman*.]
          \item 证明对某些 $g_{1},g_{2}\in\mathbb{R}[x]$, 有
            \[
              f-g_{1}-yg_{2}\in\langle x^{2}+y^{2}-1\rangle.
            \]
          \item 利用 $f\in\mathcal{I}(X)$ 证明对所有 $a\in[-1,1]$, 都有
            \[
              g_{1}(a)^{2}=(1-a^{2})g_{2}(a)^{2},
            \]
            并得出 $g_{1}(x)^{2}=(1-x^{2})g_{2}(x)^{2}$.
          \item 利用分解的唯一性证明 $g_{1}$ 和 $g_{2}$ 都是零多项式, 由此得到 $f\in\langle x^{2}+y^{2}-1\rangle$.
        \end{enumerate}
    \end{enumerate}
  \item\label{ex:vanishing-ideal-of-unit-n-sphere} 设 $X=\mathcal{V}(x_{1}^{2}+\cdots+x_{n}^{2}-1)\subseteq\mathbb{A}^{n}_{\mathbb{R}}$ 是单位 $n$ 维球面. 推广上一道习题, 证明 $\mathcal{I}(X)=\langle x_{1}^{2}+\cdots+x_{n}^{2}-1\rangle$.
  \item\label{ex:unit-circle-infinite-field} 设 $K$ 是满足 $1\ne-1$ 的无限域, 且 $X=\mathcal{V}(x^{2}+y^{2}-1)\subseteq\mathbb{A}^{2}_{K}$.
    \begin{enumerate}[label=(\alph*)]
      \item 对任意满足 $c^{2}\ne-1$ 的 $c\in K$, 证明
        \[
          \left(\frac{c^{2}-1}{c^{2}+1},\frac{2c}{c^{2}+1}\right)\in X.
        \]
      \item 证明有无穷多个 $a\in K$, 使 $(a,b)\in X$ 对某个 $b\in K$ 成立.
      \item 证明 $\mathcal{I}(X)=\langle x^{2}+y^{2}-1\rangle$, 可以改用 Exercise~\ref{ex:vanishing-ideal-of-real-unit-circle} 的策略来处理本题.
    \end{enumerate}
  \item\label{ex:parabola-vanishing-ideal-over-f2} 设 $K=\mathbb{F}_{2}$ 是含两个元素的有限域. 令
    \[
      X=\mathcal{V}(y-x^{2})\subseteq\mathbb{A}^{2}_{\mathbb{F}_{2}}.
    \]
    证明 $\mathcal{I}(X)\ne\langle y-x^{2}\rangle$, 这与 Example~\ref{ex:vanishing-ideal-of-parabola} 中讨论的 $\mathbb{R}$ 上的类似情形不同. 求 $\mathcal{I}(X)$.
  \item\label{ex:smallest-affine-variety-containing-set} 设 $X\subseteq\mathbb{A}^{n}$ 是任意子集.
    \begin{enumerate}[label=(\alph*)]
      \item 证明 $\mathcal{V}(\mathcal{I}(X))$ 是包含 $X$ 的最小仿射簇, 其含义如下: 若 $Y\subseteq\mathbb{A}^{n}$ 是任意仿射簇, 且 $X\subseteq Y$, 则 $\mathcal{V}(\mathcal{I}(X))\subseteq Y$.
      \item 任选一个不是仿射簇的 $X\subseteq\mathbb{A}^{2}_{\mathbb{R}}$, 计算 $\mathcal{V}(\mathcal{I}(X))$, 以说明 (a).
    \end{enumerate}
\end{BookExercises}
